\subsection{Clifordo algebros}


Toliau nagrinėsime nagrinėsime dvi Clifordo algebras $\mathup{Cl}^+(3)$ ir
$\mathup{Cl}^-(3)$. Tai yra vectorinės erdvės su papildoma bitiesine daugybos
operacija, kuri leidžia sudauginti du algebros vektorius ir gauti naują
vektorių algebroje. Šias algebras apibrėšime pirma paimdami daug bendresnę
tenzorinę algebrą ir tada jai sutekdami norimą struktūrą.

\begin{definition}
	Vienetinė tenzorinė algebra $T(\set R^3)$ yra tiesioginė suma
	\begin{equation}
		T(\mathbb{R}^3) 
		=
		\bigoplus_{k=0}^{\infty} T^k
		=
		\mathbb{R} 
		\oplus \set R^3 
		\oplus (\set R^3 \otimes \set R^3) 
		\oplus (\set R^3 \otimes \set R^3 \otimes \set R^3) 
		\oplus \dots
		\,,
	\end{equation}
	$T^k$ čia yra vektorinės erdvės $T^0 = \set R$\: ,\: $T^1=\set R^3$\: ,\: $T^2=\set R^3 \otimes \set R^3$\: ,\:
	$T^3=\set R^3 \otimes \set R^3 \otimes \set R^3$
	ir taip toliau. Algebros daugyba yra duodama vektorių tenzorinės sandaugos.
	Vienetinis tenzorinės algebros elementas yra žymimas $\mathbb 1$\:.
\end{definition}
Kitaip paskius tai yra algebra, kurios vektoriai yra ivairių ilgiu tenzorinių
sandaugos $\vb v_1 \otimes \vb v_2 \otimes \dots \otimes \vb v_n$\:,\: $\vb v_i \in \set R^3$
ir tokių sandaugų sumos.
\begin{example}
	Vektorius $3\cdot\mathbb 1 + \vb v$ priklauso $T(\set R^3)$ ir yra suma vektorių iš $T^0$ ir $T^1$.
	Vektorius $2\vb w + \vb v \otimes \vb w$ yra suma vektorių iš $T^1$ ir $T^2$.
	Jų sandauga tenzorinėje algebroje yra
	\begin{equation}
		\begin{split}
			\qty(3\cdot\mathbb 1 + \vb v)
			\otimes
			\qty(2\vb w + \vb v \otimes \vb w)
			=&
			3\cdot\mathbb 1 \otimes 2\vb w
			+ 3\cdot\mathbb 1 \otimes \vb v \otimes \vb w
			+ \vb v \otimes 2\vb w
			+ \vb v \otimes \vb v \otimes \vb w
			\\=&
			3\cdot2\qty(\mathbb 1 \otimes \vb w)
			+ 3\qty(\mathbb 1 \otimes \vb v \otimes \vb w)
			+ 2(\vb v \otimes \vb w)
			+ \qty(\vb v \otimes \vb v \otimes \vb w)
			\\=&
			6 \vb w
			+ 3\qty(\vb v \otimes \vb w)
			+ 2(\vb v \otimes \vb w)
			+ \qty(\vb v \otimes \vb v \otimes \vb w)
			\\=&
			6 \vb w
			+ 5\qty(\vb v \otimes \vb w)
			+ \qty(\vb v \otimes \vb v \otimes \vb w)
			\,,
		\end{split}
	\end{equation}
	čia pirmoje eilutėje pasinaudojome tuo kad tenzorinės sandauga
	distributyvi; antroje eilutėje naudojames tuo kad tenzorinė sandauga
	tiesiška kiekvieno dauginamo vektoriaus atžvilgiu, kad iškeltųme 
	skaliarinius daugiklius; trečioje eilutėje sudedame tiesiškai
	priklausomus vektorius.
	Gautas vektorius yra suma vektorių iš $T^1$, $T^2$ ir $T^3$.
\end{example}


Dabar $T(\set R^3)$ suteiksime papildomą strutūrą, kad algebroje
$\mathup{Cl}^+(3)$ būtų tenkinamas sąryšis
\begin{equation}
	\vb v \otimes \vb w
	+
	\vb w \otimes \vb v
	=
	2(\vb v , \vb w) \mathbb 1
	\quad ,
	\label{eq:cl+}
\end{equation}
o algebroje $\mathup{Cl}^+(3)$ tenkinamas sąryšis
\begin{equation}
	\vb v \otimes \vb w
	+
	\vb w \otimes \vb v
	=
	- 2(\vb v , \vb w) \mathbb 1
	\quad .
	\label{eq:cl-}
\end{equation}
Šiuos sąryšius ,,užkoduojame'' į aibes $S^\pm \subset T(\set R^3)$
\begin{equation}
	\begin{split}
		S^+ = \qty{
			\vb v \otimes \vb w
			+
			\vb w \otimes \vb v
			-
			2(\vb v , \vb w) \mathbb 1
			\mid
			\vb v, \vb w \in \set R^3
		}
		\\
		S^- = \qty{
			\vb v \otimes \vb w
			+
			\vb w \otimes \vb v
			+
			2(\vb v , \vb w) \mathbb 1
			\mid
			\vb v, \vb w \in \set R^3
		}\,.
	\end{split}
\end{equation}
Užrašome tenzorinės algebros poalgebrius $I^\pm \subset T(\set R^3)$,
gauamus iš visų vektorių, kurie turi bent vieną daugiklį iš $S^\pm$
\begin{equation}
	I^+ = \qty{
		\sum_{\substack{
			s \in S^\pm \\ 
			a,b \in \mathup{T(V)} 
		}}
		a \otimes s \otimes b
	}
	\quad .
\end{equation}
Šie poalbebriai papildomai tenkina savybę, kad betkokiem $a \in T(\set R^3)$
ir $d \in I^\pm$, $a\otimes d, d \otimes a \in I^\pm$. Poalgebris
tenkinantis šia savybę yra vadinamas algebros idealu.
Literatūroje šią konstrukciją yra įprasta pateikti sakiniu:
,,$I^\pm$ yra algebros $T(\set R^3)$ idealas genaruojams 
sąryšio (\ref{eq:cl+}) ar atitinkamai (\ref{eq:cl-})''.

\begin{definition}
	Algebra $\mathup{Cl}^+(3)$ yra faktorinė algebra $T(\set R^3)/I^+$
	ir algebra $\mathup{Cl}^-(3)$ yra faktorinė algebra $T(\set R^3)/I^-$.
\end{definition}
Faktorinės algebros elementai yra šoninės klasės sudarytos iš elementų,
kuriuos galime gauti vieną iš kito pridedant vektorių iš idealo.
\begin{example}
	$2\vb{(v,w)} \mathbb 1 - \vb v \otimes \vb w + I^+$ yra šoninė klasė,
	kuri yra vektorius algebroje $\mathup{Cl}^+(3)$. 
	$2\vb{(v,w)} \mathbb 1 - \vb v \otimes \vb w$ 
	yra šoninės klasės atstovas, o ,,$+ I^+$'' žymi
	kad šiai klasei taip pat prikauso visi vectoriai gaunami prie
	klasės atstovo pridėjus vektorių iš $I^+$.
	Tai mums leidžia atlikti skaičiavimus kaip
	\begin{equation}
		\vb{(v,w)} \mathbb 1 - \vb v \otimes \vb w + I^+
		= 
		\vb{(v,w)} \mathbb 1 - \vb v \otimes \vb w 
		+ \qty(
			\vb v \otimes \vb w + \vb w \otimes \vb v 
			- 2\vb{(v,w)}\mathbb 1
		) 
		+ I^+
		=
		\vb w \otimes \vb v + I^+
		\,,
	\end{equation}
	čia pasinaudojome tuo kad 
	$\qty( \vb v \otimes \vb w + \vb w \otimes \vb v - 2\vb{(v,w)}\mathbb 1) \in I^+$.
	Pavyko parodyti, kad $2\vb{(v,w)} \mathbb 1 - \vb v \otimes \vb w$ 
	ir $\vb w \otimes \vb v$ atstovaują tą pačią klasę.
\end{example}

\begin{example}
	Tegul $a,b \in T(\set R^3)$ ir
	$
		a \otimes \vb{(v \otimes w + v \otimes w)} \otimes b + I^-
	$
	yra šoninė klasė, kuri yra vektorius iš $\mathup{Cl}^-(3)$.
	Šią išraišką supapranstiname atlidami veiksmus
	\begin{equation}
		\begin{split}
			a \otimes \vb{(v \otimes w + v \otimes w)} \otimes b + I^-
			=&
			a \otimes \vb{(v \otimes w + v \otimes w)} \otimes b 
			- a \otimes \qty(\vb{v \otimes w + v \otimes w}+2\vb{(v,w)} \mathbb 1)  \otimes b 
			+ I^-
			\\=&
			a \otimes \qty(
				\vb{(v \otimes w + v \otimes w)}
				-\qty(\vb{v \otimes w + v \otimes w}
				-2\vb{(v,w)} \mathbb 1) 
			)\otimes b 
			+ I^-
			\\=&
			a \otimes \qty(
				-2\vb{(v,w)} \mathbb 1 
			)\otimes b 
			+ I^-
			\\=&
			-2\vb{(v,w)} \qty(a \otimes \mathbb 1 \otimes b)
			+ I^-
			\\=&
			-2\vb{(v,w)} \qty(a \otimes b)
			+ I^-
			\,.
		\end{split}
	\end{equation}
	Parodėme, kad sudarius faktorinę algebrą $\mathbb{Cl}^-(3)$,
	galime jos viduje taikyti sąryšį (\ref{eq:cl-})
	manipuliuojant šoninės klasės atstovus.
\end{example}
\begin{example}
	Tegul $a$ ir $b$ yra šoninių klasių atstovai iš
	$\mathup{Cl}^\pm(3)$.
	Jų sandauga algebroje yra
	\begin{equation}
		(a+I^\pm)\otimes(b+I^\pm)
		=
		a \otimes b
		+ a \otimes I^\pm
		+ I^\pm \otimes b
		+ I^\pm \otimes I^\pm
		=
		a \otimes b + I^\pm
		\,.
	\end{equation}
	Sumos nariai $
		a \otimes I^\pm
		= I^\pm \otimes b
		= I^\pm \otimes I^\pm
		= I^\pm
	$, nes $I^\pm$ yra algebros $T(\set R^3)$ idealas.
	Tai
	\begin{equation}
		(a+I^\pm)\otimes(b+I^\pm)
		=
		a \otimes b + I^\pm
		\,.
	\end{equation}
\end{example}
Toliau priimame sutrumpintą žymėjimą, nes išreikštai
su šoninėmis klasėmis nebedirbsime. Užrašysime tik
kalsės atstovą, o kad išraiška skirtūsi nuo 
vektorių tenzorinės sandaugos vietoje ženklo ,,$\otimes$''
tiesiog rašysime vektorius vieną šalia kito.
Vietoje
\begin{equation}
	a \otimes \vb{(v \otimes w + v \otimes w)} \otimes b + I^-
	=
	-2\vb{(v,w)} \qty(a \otimes b) + I^-
\end{equation}
rašysime
\begin{equation}
	a \vb{(v w + v w)} b
	=
	-2\vb{(v,w)} a b
	\,.
\end{equation}

Štai kelios svarbios tapatybės
\begin{preposition}
	\begin{equation}
		\vb v \vb v 
		= \frac{1}{2}(\vb v \vb v + \vb v \vb v) 
		= \pm(\vb v,\vb v) \mathbb 1
		\,.
	\end{equation}
\end{preposition}

\begin{preposition}
	Algebroje $\mathup{Cl}^+(3)$ 
	\begin{equation}
		-\vb v \vb w \vb v 
		= -\qty( 2(\vb v,\vb w)\mathbb 1 - \vb{wv} )\vb v 
		= -2\vb{(v,w)} \vb v + \vb{wvv} 
		= -2\vb{(v,w)} \vb v + \vb{(v,v)}\vb w 
		\,.
	\end{equation}
\end{preposition}

\begin{preposition}
	Algebroje $\mathup{Cl}^-(3)$ 
	\begin{equation}
		\vb v \vb w \vb v 
		= \qty( - 2(\vb v,\vb w)\mathbb 1 - \vb{wv} )\vb v 
		= -2\vb{(v,w)} \vb v - \vb{wvv} 
		=  -2\vb{(v,w)} \vb v + \vb{(v,v)}\vb w 
		\,.
	\end{equation}
\end{preposition}

Atkreipiame dėmesį, kad kai $(v,v)=1$,
$\vb{(v,w)v}$ yra vektoriaus $w$ projekcija į $v$
ir $-2\v{(v,w)v} + \vb w$ yra vektoriaus atspindys
per plokštumą statmeną vektoriui $v$.

Tegul $\qty{\vb e_1,\vb e_2,\vb e_3}$ yra ortonormali $\set R^3$ bazė. Bazinius vektorius galime perketi
į algebra $\mathup{Cl^\pm(V)}$, kur jie turi tenkinti
\begin{equation}
	\begin{split}
		\vb e_i \vb e_j &= - \vb e_j \vb e_i \ , \ i \neq j \\
		\vb e_i \vb e_i &= \pm \mathbb 1
		\quad .
	\end{split}
\end{equation}
Su šių sąryšių pagalba galime bet kokį $\mathup{Cl^\pm(V)}$ elementą užrašyti, kaip
tiesinę kombinaciją aštuonių elementų:
\begin{equation}
		\mathbb 1 ,\,
		\vb e_1 ,\, \vb e_2 ,\, \vb e_3 ,\, 
		\vb e_2 \vb e_3 ,\, \vb e_1 \vb e_3 ,\, \vb e_1 \vb e_2 ,\,
		\vb e_1 \vb e_2 \vb e_3 ,
\end{equation}
nes su pirma taisykle betkurį tenzorinės algebros bazės elementą galima išrikiuoti taip kad
ortų indeksai būtų surašyti didėjimo tvarka ir su antra taisykle pasikartojantys ortai panaikinami.
Tai reiškia $\mathup{Cl^\pm(V)}$ yra aštuonmatės algebros. Algebros elemento norma $\abs{v}$ yra
abibrėžiama, taip pat kaip vektoriaus ilgis
\begin{equation}
	\abs{v} = \sqrt{ \sum_{I \subset {1,2,3}} {v_I}^2 }
	\quad ,
\end{equation}
čia summa eina per visus indeksų poaibius, kurie atitinka bazinius algebgros elementus.
Bendresniu atvėju $\mathup{dim}\qty[\mathup{Cl}^\pm\qty(\set R^n)] = 2^n$.

Apibrėžiame "transponuotą elementą" apibėždami kaip transponavimas veikia ant algebros bazės
\begin{equation}
	\qty(\vb e_{i_1} \vb e_{i_2} \dots \vb e_{i_k} )^t
	=
	\vb e_{i_k} \vb e_{i_{k-1}} \dots \vb e_{i_1}
	=
	(-1)^{k-1} \vb e_{i_1} \vb e_{i_2} \dots \vb e_{i_k} 
	\quad .
\end{equation}
Transpozicija yra algebros anti-homomorfizmas, tai reiškia, kad betkuriems algebros elementams $v$ ir $w$ galioja
\begin{equation}
	(v w)^t = w^t v^t
	\quad .
\end{equation}


Taip pat veikimu ant bazės apibrėžiame algebros homomorfizmą
\begin{equation}
	\alpha( \vb e_{i_1} \vb e_{i_2} \dots \vb e_{i_k} )
	= 
	(-1)^{k} \vb e_{i_1} \vb e_{i_2} \dots \vb e_{i_k} 
	\quad ,
\end{equation}
jis algebros elemento komponentes prie lyginio ortų skaičiau siunčia į savo minusą, o
komponentčių prie nelyginio skaičiaus nepakeičia.











